
\section{Kredyty mieszkaniowe}
\label{sec:kredyty}

Pan A. kupił mieszkanie w Warszawie za X zł. Część kwoty zakupu musiał finansować kredytem hipotecznym, na kwotę K zł.
Od dnia d1 zaczął spłacać kredyty hipoteczny.
Do dnia d2 spłacił Y zł kapitału kredytu hipotecznego. Ciąg rat wynosi R = [r1, r2, r3, ...].
W dniu d2 Pan A. sprzedał mieszkanie za kwotę S zł.


Zamiast kupować mieszkanie Pan A. mógłby zainwestować pieniądze w fundusz inwestycyjny, który przynosiłby mu zyski na poziomie Z\% rocznie.

Oblicz za ile Pan A. powinien sprzedać mieszkanie, żeby opłacało mu się kupować mieszkanie na kredyt hipoteczny, a nie inwestować pieniądze w fundusz inwestycyjny.